\subsection{\texorpdfstring{DERIVATIVE OF \(\log_a x\)}{DERIVATIVE OF LOG BASE A OF X}}

Since \(a^x\) is strictly monotone on \(\mathbf{R}\) for \(a \neq 1\) , applying the rule for differentiating inverse functions (I, p.~17, prop. 6) gives, for all \(x > 0\)

\[
\mathrm{D} \big (\log_ {a} x \big) = \frac {1}{x \log a}\tag{9}
\]

and in particular

\[
\mathrm{D} (\log x) = \frac {1}{x}.\tag{10}
\]

If u is a real function admitting a derivative at the point \(x_{0}\) and such that \(u(x_{0}) > 0\) , then the function \(\log u\) admits a derivative equal to \(u'(x_{0})/u(x_{0})\) at the point \(x_{0}\) . In particular, we have \(\mathrm{D}\left(\log |x|\right) = 1/|x| = 1/x\) if x \textgreater{} 0, and

\[
\mathrm{D} \big (\log | x | \big) = - \frac {1}{| x |} = \frac {1}{x}
\]

if x \textless{} 0; in other words, \(\mathrm{D}\big(\log|x|\big) = 1/x\) for any \(x \neq 0\) . One concludes that if, on an interval I, the real function u is not zero and admits a finite derivative, then \(\log|u(x)|\) admits a derivative equal to \(u'/u\) on I; this derivative is called the logarithmic derivative of u. It is clear that the logarithmic derivative of \(|u|^{\alpha}\) is \(\alpha u'/u\) , and that the logarithmic derivative of a product is equal to the sum of the logarithmic derivatives of the factors; these rules often provide the fastest way to calculate the derivative of a function. They give again, in particular, the formula

\[
\mathrm{D} \left(x ^ {\alpha}\right) = \alpha x ^ {\alpha - 1}
\]

\[
(\alpha \text {   an   arbitrary   real   number,   } x > 0)\tag{11}
\]

which has already been shown by another method (II, p.~69).

Example. If u is a function \(\neq 0\) on an interval I, and v is any real function, then \(\log(|u|^{v}) = v \cdot \log |u|\) , so if u and v are differentiable

\[
\frac {1}{| u | ^ {v}} \mathrm{D} (| u | ^ {v}) = v ^ {\prime} \log | u | + v \frac {u ^ {\prime}}{u}.
\]

